\section{Prose documents}\label{sec:prose-documents}

Syllabi, TA information sheets, course policies -- continuous prose rather
than problems -- are \file{CourseDocument} documents. It is the ``designed''
class of the three: coloured headings, a titled cover, a table of contents,
and a text measure tuned for reading. (This manual is itself a
CourseDocument.) The demo repo's \file{_syllabus/syllabus.tex}:

\begin{manexsource}
%! views: instructor
\documentclass[11pt]{CourseDocument}

\title{\course{101} --- Syllabus}
\subtitle{%
  Fall 2026, University of New Brunswick \\[1ex]
  \normalfont\normalsize\itshape Sections FR01 and FR02.
}

\begin{document}

\maketitle
\pagestyle{toc}
\tableofcontents

\newpage
\pagestyle{myfancy}

In \course[MATH]{1013} we will be investigating...

\end{document}
\end{manexsource}

That opening sequence is the documented idiom, worth copying verbatim.
\cs{maketitle} runs while the page style is still \texttt{plain}, so the
title page stays bare; the contents page then opts into \texttt{toc}
(footer page number only); and \texttt{myfancy} -- a small-caps running
head with the current section, page number in the footer -- starts with the
body. Neither named style activates itself: each document chooses when.

The pieces in play:

\begin{itemize}
  \item \cs{title} and \cs{subtitle} feed a shared title block that also
    stamps \emph{Updated \textnormal{(date, time)}} -- these documents are
    living and revised mid-term, and the stamp says which printing is
    current. The subtitle's second line above shows the usual trick for a
    quieter sub-subtitle: reset the font inline.
  \item \cs{course}\texttt{[MATH]\{1013\}} sets a course designation in the
    body serif with lining numerals, tied so it cannot break across lines --
    consistent wherever it appears, prose or title. The bracket defaults to
    the course's own prefix, so \cs{course}\texttt{\{101\}} works too.
  \item Headings, the title block and internal links all carry the course
    identity colour -- named \emph{once} (\env{CourseIdentity}, a mahogany)
    and used unmodified, so restyling a course document is a one-line
    change. Deliberately \emph{not} green: green means ``solutions, not for
    students'' everywhere else in the toolkit, so a student-facing document
    must not read as green.
  \item \texttt{[legalpaper]} is available, as on assessments; stock
    \file{article} options pass through.
\end{itemize}

Body text in this class uses old-style figures -- the numerals with
ascenders and descenders you see in this manual's prose -- because these
documents are continuous reading. Dates, office numbers and the like sit
quietly in the text; \cs{liningnums}\texttt{\{2026\}} forces lining figures
where a number should read as data, and math and \cs{course} numerals stay
lining automatically. The full figure policy lives with
\file{Typefaces.sty} (section~\ref{sec:ref-typefaces}).

Not every prose document wants them, though. A syllabus is only half
continuous reading: the other half is \emph{data} scanned rather than
read -- test dates, room numbers, section codes, percentage weights --
and a page of those in old-style figures reads as restless. When
\cs{liningnums} starts appearing more often than not, stop reaching for
it and set the class option instead:

\begin{manexsource}
\documentclass[11pt,lining]{CourseDocument}
\end{manexsource}

\noindent It is a class option and not something a preamble can do for
itself, because Typefaces is loaded during \cs{documentclass} -- by the
time a document could pass it an option, it is already in. The two
choices are exclusive: name \texttt{lining} and the class passes exactly
that to Typefaces in place of \texttt{osf}.

Two things CourseDocument deliberately lacks: the exercise system (a
syllabus imports no problems -- the class loads no Exercises/Workspace
machinery), and any use of \cs{workspace}. It \emph{does} carry the
instructor-notes layer, so a syllabus can hold staff-only annotations --
office-hour strategy, grading-scheme rationale -- and ship an
\file{-INSTRUCTOR} view like any other document
(section~\ref{sec:build-and-find}).

\subsection{Passages}\label{sec:ch5-passages}

A \textbf{passage} is prose that belongs to the author rather than to the
course: why the subject is worth the effort, how marking works, what makes a
solution readable. Because it is shared across courses it lives in its own
content repository, mounted beside the topic submodules and reached by
relative path --- it is content, not a tier (see the architecture chapter).

The reason this is a command rather than a bare \cs{input} is \textbf{heading
level}. The same passage is a \cs{section} of one document and a
\cs{subsection} of another, and a passage carrying headings of its own has to
move with its title. So the file is written at its \emph{natural} depth ---
\cs{PassageTitle} for its own title, plain \cs{subsection} for anything
inside it --- and the importing document says where that lands:

\begin{manexsource}
% teaching/passages/why-learn-math.tex
\PassageMeta{
  title = {Why Learn Math?},
  label = {sec:whylearnmath},
}
\PassageTitle
Mathematics is beautiful, powerful, and useful.
\end{manexsource}

\begin{manexsource}
\importpassage{../teaching/passages/}{why-learn-math.tex}
\importpassage[level=subsection]{../teaching/passages/}{why-learn-math.tex}
\end{manexsource}

\noindent \cs{PassageTitle} takes \emph{no arguments}. Its title and its
\cs{label} both come from the \cs{PassageMeta} block above it --- the same block
\file{find-passages} reads off disk (section~\ref{sec:ref-passagemeta}). The
label names the passage's own idea, so it belongs to the passage rather than to
whichever document imports it.

Unlike \cs{ProblemMeta} and \cs{FactMeta}, which are validated and thrown away,
this block is read at build time, so it is not optional: a passage with no
\env{title} is an error rather than an empty heading.

The shift is three \cs{let}s inside a group, run top-down so each captures the
still-original meaning of the level below. \cs{PassageTitle} then needs no
special case --- it is \cs{section}, already pointed at the right target ---
and \cs{label}, \cs{ref}, starred forms and table-of-contents entries keep
working because these are the real sectioning commands.

\textbf{One level of headroom, and deeper is an error.} This class styles
three heading levels and \cs{paragraph} is not among them, so \texttt{level}
accepts \texttt{section} (the default) or \texttt{subsection} and refuses
anything further rather than degrading silently. A passage may therefore carry
one rank of headings below its title and no more. In practice that is the
constraint that decides how a passage is cut: a body of prose with three
ranks is two passages, not one.

A passage needing a number the course owns names a macro --- say
\cs{PointsPerQuestion} --- and lets the course define it in its own
\file{courseinfo.sty}. There is no machinery for that on purpose: a course
that imports a passage without defining what it needs gets an
undefined-control-sequence error, which is the loud failure wanted. Writing
\emph{around} a parameter travels further than introducing one.
